\section{Proposed Work}
\label{sec:proposed_work}

Based on the literature reviewed in Section~\ref{sec:literature_review}, this project's radiomics pipeline should be built as follows: FDG-PET input taken from ADNI's own harmonized \texttt{UCBERKELEYFDG} processing output rather than from re-coregistered raw DICOM/ECAT/Interfile images \cite{jagust2024petcore,jagust2010petcore}; regions of interest drawn from ADNI's existing FreeSurfer parcellation tables (\texttt{UCSFFSX51}) with explicit attention to whether posterior cingulate, occipital, and parietal regions are separable at sufficient granularity, given their established role in the DLB cingulate island sign \cite{mckeith2017dlb,lim2009cingulateisland}; feature extraction via PyRadiomics using the IBSI-standard feature-class set with wavelet filtering \cite{vangriethuysen2017pyradiomics,zwanenburg2020ibsi}; and a classifier stage built around nested cross-validation, with feature selection and class-balancing performed strictly inside each training fold \cite{demircioglu2021bias,demircioglu2024oversampling}.

\subsection{A required correction: decay-corrected frame combination}

A concrete, previously unaddressed correction surfaced by this review concerns the handling of dynamic FDG-PET acquisitions. A subset of this project's cohort -- 37 series across 25 subjects, all early-phase, HRRT-scanner acquisitions with no fallback DICOM/ECAT visit -- exist only in Interfile format and required a custom reader (\texttt{src/dlb\_radiomics/interfile.py}), which currently combines each series' six dynamic frames by a plain raw-count sum. Standard PET reconstruction and processing practice, as documented in ADNI's own PET Core protocol \cite{jagust2010petcore}, decay-corrects each frame to a common reference time (typically the tracer injection time) before combining frames into a static or standardized-uptake-value image; a plain sum of raw, non-decay-corrected frame counts systematically under-weights later frames relative to earlier ones as the tracer decays, biasing the resulting static image in a way that would not be present in the DICOM- and ECAT-derived majority of the cohort. This is not a stylistic preference but a methodological gap that should be closed before these 37 series are used for radiomic feature extraction: the Interfile reader (and the equivalent multi-frame handling for the 179 ECAT7 series) should apply per-frame decay correction, using each frame's start time and the F-18 physical half-life (109.77 minutes), before summing.

\subsection{Remaining scope not addressed by this review}

This review did not attempt to resolve the directory layout for \texttt{src/}, since that follows mechanically once the coregistration and ROI-source decisions above are fixed, nor did it find literature specific enough to override the project's existing default PyRadiomics feature-class configuration -- no DLB- or FDG-PET-specific study was found recommending a narrower or different feature-class subset than the IBSI default set. Both remain open implementation tasks rather than open research questions.
